% Generated draft: requires mathematical and source review.
\documentclass[11pt]{article}
\usepackage[margin=25mm]{geometry}
\usepackage{fontspec}
\setmainfont{Latin Modern Roman}
\usepackage{amsmath,amssymb,mathtools,unicode-math}
\setmathfont{Latin Modern Math}
\usepackage{enumitem,needspace,xurl,hyperref}
\setlength{\parindent}{0pt}
\setlength{\parskip}{5pt}
\setlength{\emergencystretch}{3em}
\newcommand{\N}{\mathbb{N}}
\newcommand{\Z}{\mathbb{Z}}
\newcommand{\Q}{\mathbb{Q}}
\newcommand{\R}{\mathbb{R}}
\newcommand{\C}{\mathbb{C}}
\newcommand{\F}{\mathbb{F}}
\newcommand{\A}{\mathbb{A}}
\newcommand{\PP}{\mathbb{P}}
\newcommand{\E}{\mathbb{E}}
\newcommand{\eps}{\varepsilon}
\newcommand{\dd}{\,\mathrm{d}}
\begin{document}
\section*{UnsolvedMath ID: 4600004}
\textbf{Source problem code:} catalogue entry
\subsection{Definitions and mathematical statement}
Write \(\Q_+=\{q\in\Q:q\ge0\}\), the nonnegative rational semiring. Positive matrices here have every entry strictly greater than zero. Let \(A\in M_m(\Q_+)\) and \(B\in M_n(\Q_+)\) be positive square matrices, where \(m,n\ge1\).

They are shift equivalent over \(\Q_+\) if there are a lag \(\ell\ge1\) and rectangular nonnegative rational matrices \(R\in M_{m\times n}(\Q_+)\), \(S\in M_{n\times m}(\Q_+)\) such that
\[
AR=RB,\qquad SA=BS,\qquad A^\ell=RS,\qquad B^\ell=SR.
\]
They are elementary strong shift equivalent when \(A=RS\) and \(B=SR\) for rectangular matrices over \(\Q_+\). They are strong shift equivalent over \(\Q_+\) when there is a finite chain
\[
A=A_0,A_1,\ldots,A_k=B
\]
of square nonnegative rational matrices, of possibly different sizes, with each consecutive pair elementary strong shift equivalent over \(\Q_+\). Intermediate matrices and factors need not have strictly positive entries.

The statement to prove is
\[
\forall m,n\ge1\ \forall A\in M_m(\Q_+),B\in M_n(\Q_+)
\text{ positive},\qquad
A\sim_{\mathrm{SE},\Q_+}B\ \Longrightarrow\ A\sim_{\mathrm{SSE},\Q_+}B,
\]
where the displayed relation symbols abbreviate exactly the definitions above. The supplied research summary records this positive rational assertion as solved by Kim--Roush (1990). It also reports real/rational variants and path results; the present target retains rational endpoints and rational nonnegative factors, without extending the conclusion to unrestricted real endpoint matrices.
\subsection{Short English statement}
Positive square rational matrices that are shift equivalent using nonnegative rational matrices should be connected by a finite chain of elementary strong shift equivalences over the same semiring. The supplied registry records this as a literature theorem.
\subsection{Sources}
\begin{itemize}
\item Local UnsolvedMath registry, permanent ID \(4600004\), exact source code \texttt{catalogue entry}, ``Positive rational shift equivalence.''
\item Boyle, \emph{Open Problems in Symbolic Dynamics} (2008), Problem/Question/Conjecture 5.1; \url{https://www.math.umd.edu/~mboyle/papers/openfinalsub3nov2007.pdf}. Recorded access 2026-07-29; pdf-text extraction.
\item Supplied research summary attributes the affirmative theorem to K. H. Kim and F. W. Roush (1990), and quotes a Boyle Kansas talk saying the positive shift-equivalence conjecture holds with \(\R\) or \(\Q\) in place of \(\Z\). Complete publication and talk details are not supplied.
\item Boyle--Kim--Roush, ``Path methods for strong shift equivalence of positive matrices,'' as supplied: the summary reports descent from real nonnegative SSE to rational nonnegative SSE for positive rational endpoints, and SSE along paths of positive shift-equivalent real matrices. No full bibliographic details are supplied.
\end{itemize}
\end{document}
